\chapter{Notation, Open Problems, and Numerical Index}
\label{app:base-backmatter}

\section{Notation}

This supplement uses the symbols of the main manuscript (\Cref{tab:notation}) and of Supplement G
(\Cref{tab:gaussian-notation}) with their meanings there. \Cref{tab:base-notation} lists the symbols
that recur across the chapters of Parts B-I and B-II; a symbol used inside one chapter is declared at
its opening, and the lattice part declares its symbols in \Cref{tab:rg-lattice-notation}. The word
pullback denotes the connection-relative pullback of fiber tensors to the base of
\Cref{def:pb-informational-pullbacks}; restriction of a tensor along an injective map of parameters,
which Supplement G also calls a pullback, is a different construction.

{\footnotesize
\renewcommand{\arraystretch}{1.12}
\begin{longtable}{@{}>{\raggedright\arraybackslash}p{0.29\textwidth}
>{\raggedright\arraybackslash}p{0.46\textwidth}
>{\raggedright\arraybackslash}p{0.19\textwidth}@{}}
\caption{Notation of Supplement B.}
\label{tab:base-notation}\\
\toprule
\textbf{Symbol} & \textbf{Meaning} & \textbf{Introduced} \\
\midrule
\endfirsthead
\toprule
\textbf{Symbol} & \textbf{Meaning} & \textbf{Introduced} \\
\midrule
\endhead
\midrule
\multicolumn{3}{r@{}}{\emph{continued on the next page}} \\
\endfoot
\bottomrule
\endlastfoot
\multicolumn{3}{@{}l}{\textbf{Base transport and multi-point designs}} \\
$f$; $s$ & a base map; a section & \cref{ch:base-introduction} \\
$u_i^b$, $u_i^m$, $h_i$ & belief and model frames of agent $i$ in two atlases; the relative frame, $u_i^m=u_i^b\cdot h_i$ & \cref{prop:geo-relative-frame} \\
$g_i^b$, $g_i^m$ & passive re-expressions of the two atlases, $u_i^{x\prime}=u_i^x\cdot(g_i^x)^{-1}$; the common frame has $g_i^b=g_i^m=g_i$ & \cref{prop:geo-relative-frame} \\
$\omega_b$, $\omega_m$; $H^b$, $H^m$; $H^x_eX$ & principal connections; their horizontal distributions; the horizontal lift of $X\in T_c\mathcal C$ at a point $e$ of $(\mathcal E_x)_c$ & \cref{def:geo-connections} \\
$\gamma$; $\Omega_\gamma$, $\widetilde\Omega_\gamma$ & piecewise-smooth base curve; parallel transports of belief and model laws along it & \cref{def:geo-connections} \\
$A_i^x=(u_i^x)^*\omega_x$ & local connection form in the frame $u_i^x$ & \cref{prop:geo-connection-frame-laws} \\
$\mathcal D\Phi$, $\mathcal D\widetilde\Phi$ & covariant defects of the cross morphisms & \cref{def:geo-covariant-defects} \\
$G_b\times G_m$ & structure group of the optional product-gauge extension & \cref{def:geo-product-extension} \\
$D=\{c_d\}$, $d$; $I_D$, $D_i$, $V_d$ & finite design and its index; agent--design incidence set, design points of agent $i$, agents at the design point $c_d$ & \cref{def:prob-finite-design} \\
$D_e=D_i\cap D_j$, $\Theta^x_{d,e}$, $g^x_{d,i}=g^x_i(c_d)$ & design points shared by the endpoints of $e:j\to i$; the link copy at $c_d$; the re-expression at an incidence & \cref{def:gen-shared-link-admissibility} \\
$\mathfrak S_i^x$ & measurable space of sections of channel $x$ for agent $i$ & \cref{hyp:prob-context-section-response} \\
$f_\ell:\mathcal C_\ell\to\mathcal C_{\ell+1}$; $\pi_\ell:\mathscr P_\ell\to\mathcal C_\ell$; $\mathcal P_\ell$, $\kappa_\ell$ & base map between level bases; level principal bundles; principal scale map covering $f_\ell$ and its group homomorphism & \cref{prop:rg-principal-scale-map-existence} \\
$D^\omega s$, $\operatorname{ver}^\omega$ & covariant first jet of a section; vertical projection of a connection & \cref{def:pb-informational-pullbacks} \\
$h^\omega_s$, $c^\omega_s$ & connection-relative Fisher and Amari--Chentsov pullbacks of a section & \cref{def:pb-informational-pullbacks} \\
$\Psi$, $\Psi_*=T^V\Psi$ & bundle morphism over a base map, and its vertical differential & \cref{sec:pb-coarse-naturality} \\
$A_\Psi=\mathcal D\Psi$; $\mathfrak A_{\mathcal P}$ & horizontal defect of $\Psi$; scale-connection defect form of a principal scale map & \cref{def:pb-horizontal-defect,def:pb-scale-connection-defect} \\
$\Delta_F^\Psi$; $\delta_\Psi$, $\mathcal X_\Psi$, $\mathcal Q_\Psi$ & vertical Fisher defect of $\Psi$; its base pullback and the signed and positive anomaly tensors & \cref{def:pb-base-comparison-tensors} \\
$\mathcal R$ & residual of the base Fisher defect cocycle & \cref{thm:pb-base-defect-cocycle} \\
$\mathcal Q_\ell$, $\mathsf G_\ell$ & configuration manifold of sections over a base at scale $\ell$, and its configuration metric & \cref{def:hist-finite-configuration-tier} \\
\multicolumn{3}{@{}l}{\textbf{Infinite networks and Wilsonian asymptotics}} \\
$\mathfrak X_\star$, $J_\ell$; $\widehat{\mathfrak R}_\ell$ & reference object and comparison isomorphisms into the scale-$\ell$ objects; the reference cocycle & \cref{def:rg-comparison-data} \\
$C_b$, $I_b$, $K_b=C_bI_b$; $\mathfrak R_b$ & structural coarse kernel, identification kernel, and rescaled kernel of a block ratio $b$; the measure-pair step & \cref{def:rg-rescaled-kernel} \\
$v_{\ell,a}$, $\lambda_{\ell,a}$ & compatible cross-scale mode line and its coefficients & \cref{def:rg-mode-line} \\
$b_k$; $\varsigma_{n\leftarrow\ell}$ & declared adjacent block ratios; cumulative log block scale & \cref{def:rg-mode-exponents} \\
$\chi_{n\leftarrow\ell}$, $\upsilon_{n\leftarrow\ell,a}$ & finite-time growth rate and scalar mode exponent & \cref{def:rg-mode-exponents} \\
$\mathfrak R^H_b$, $\mathfrak b^H_b$ & reference-dependent action map and discrete action beta & \cref{def:rg-action-beta} \\
$\mathfrak b^{\mathrm{ex}}_\ell$, $\mathfrak b^{\mathrm{ret}}_\ell$ & exact and retained discrete interaction beta & \cref{def:rg-interaction-beta,prop:rg-retained-beta-residual} \\
$\nabla^{\mathrm{scale}}$, $\mathfrak b^{\mathrm{scale}}$ & scale connection and continuous beta & \cref{def:rg-scale-connection} \\
$r_{\mathrm{ess}}$ & essential spectral radius, the spectral radius in the Calkin algebra & \cref{thm:rg-unital-essential-spectrum} \\
\end{longtable}
}

\section{Open problems}

\Cref{tab:base-open} lists the open problems of this supplement, each with the statement at which it
arises and a condition that would settle it. The first arises from a claim declined in
\Cref{ch:pullback-geometry} and is stated here. Open problems of the main manuscript that these
chapters extend, such as the dynamic closure of the pointwise parent
(\Cref{open:rg-pointwise-parent-dynamics}) and the slow-manifold reduction
(\Cref{open:obs-slow-manifold}), are listed there.

\openproblemheading{A canonical nondegenerate pullback geometry}{open:pb-canonical-pullback-geometry}
The pullback tensors of a fixed section are passively gauge covariant for a chosen connection
(\Cref{thm:pb-pullback-gauge-invariance}) and change by exact formulas when the connection changes
(\Cref{prop:pb-pullback-connection-change}), and where the covariant first jet has constant rank they
induce a positive-definite metric on a quotient vector bundle (\Cref{thm:pb-pullback-rank-quotient}).
The collective rule of \cref{prop:pb-collective-minimizer} selects a minimizing connection class
after common-target comparisons, law families, and weights are supplied; its isotropy component
remains undetermined. The rule does not select those additional data from the model, and a combined
belief--model tensor still needs channel weights (\Cref{hyp:pb-weighted-product-geometry}).
What is open is a pullback geometry determined by the data
of the model. It is settled by a canonical or empirically specified connection together with channel
weights under which the weighted pullback has constant rank, an integrable radical, a regular Hausdorff
leaf space, and basic induced tensors, so that
\Cref{thm:pb-pullback-rank-quotient,prop:pb-quotient-manifold-requirements} give a nondegenerate
geometry on a quotient manifold; or by a proof that no connection in a stated class yields these
properties.

{\footnotesize
\renewcommand{\arraystretch}{1.12}
\begin{longtable}{@{}>{\raggedright\arraybackslash}p{0.25\textwidth}
>{\raggedright\arraybackslash}p{0.17\textwidth}
>{\raggedright\arraybackslash}p{0.52\textwidth}@{}}
\caption{Open problems of Supplement B.}
\label{tab:base-open}\\
\toprule
\textbf{Item} & \textbf{Locator} & \textbf{Settling condition} \\
\midrule
\endfirsthead
\toprule
\textbf{Item} & \textbf{Locator} & \textbf{Settling condition} \\
\midrule
\endhead
\midrule
\multicolumn{3}{r@{}}{\emph{continued on the next page}} \\
\endfoot
\bottomrule
\endlastfoot
\multicolumn{3}{@{}l}{\textbf{Part B-I: base transport and multi-point designs}} \\
Continuum law theory & \cref{open:prob-continuum-theory} & A refining family of designs, a projective system of normalized finite laws with a Kolmogorov-type extension, a topology and sigma-algebra on a space of sections with tightness, convergence of the evidence bound, connections, and observables, control of gauges and reference measures under refinement, and a normalizable continuum interaction energy. \\
Regular frame-coordinate quotient & \cref{open:gen-gauge-quotient} & The generic stabilizers of the passive $G\times G$ coordinate action on the parameter space, and a proof of freeness and properness on declared strata on which its orbit space is a regular quotient manifold. \\
A canonical nondegenerate pullback geometry & \cref{open:pb-canonical-pullback-geometry} & A canonical or empirically specified connection with declared channel weights under which the weighted pullback has constant rank, an integrable radical, a regular Hausdorff leaf space, and basic induced tensors; or a proof that no connection in a stated class yields them. \\
Collective interface preservation & \cref{open:pb-collective-interface-preservation} & Structural conditions for coupled Gaussian quadratic-record inference and reducing linear block channels that preserve the comparisons and give a uniform positive aggregate-soldering singular-value bound through a positive horizon and nonzero coarse depth, with constants controlled along connected graphs of unbounded size. \\
Admission of declared recognition families to a configuration tier & \cref{open:hist-section-tier-admission} & For a declared family, a manifold calculus, its statistical metric, a smooth quotient, and a metric-orthogonal splitting at every scale where a natural-gradient field or a Fisher duration is asserted; or a family for which one of them provably fails while the history theorems are still invoked. \\
Smooth structure of the projectable set & \cref{open:hist-projectable-set} & A transversality or elliptic-regularity hypothesis on the base map and the recognition family under which the projectable set is proved to be a smooth submanifold carrying a smooth induced configuration map; or a family in which it provably is not. \\
Gluing of independently supplied local data & \cref{open:rg-cover-local-gluing} & A gluing theorem with hypotheses on the local laws, channels, and presentations of a cover alone, giving one joint extension compatible with every restriction and selected conditional version, or an obstruction computed from those data. \\
Graph-to-base identification & \cref{open:rg-cover-graph-base-identification} & A choice of vertex contexts, an embedding, and oriented edge curves for given graph data and a connection, with a proof that the independently specified represented links equal the corresponding base transports and hypotheses under which the realization is unique up to gauge; or graph data for which no such choice exists. \\
Connection recovery and record-class sensitivity & \cref{open:rg-cover-holonomy-recovery} & A recovery theorem for a connection or its holonomy class from a declared class of population records; negatively, a theorem that every statistic in a declared admissible controlled class is insensitive to the holonomy class. \\
Quotients beyond the chosen metric splitting & \cref{open:rg-cover-quotient-extension} & Constant-rank support data, compatible represented actions, and an admitted quotient or bundle map for which a canonical effective-support quotient, an intertwiner into a prescribed associated representation, and quotient geometry across rank changes are constructed; or an obstruction to them. \\
Informative base aggregation & \cref{open:rg-cover-base-aggregation} & A construction of informative aggregation of laws, channels, and connections along a base map that collapses contexts, with the base prior, joint base-fiber model, and observation weighting it needs, and its descent and free-energy identities; or an obstruction for a declared base map. \\
\multicolumn{3}{@{}l}{\textbf{Part B-II: infinite networks and Wilsonian asymptotics}} \\
Compatible rescaling kernels on agent networks & \cref{open:rg-rescaling-compatibility} & A family of identification kernels acting on the gauge-retained network state of the main manuscript, with a proof of the compatibility law on a declared class of finite networks at $c_*$; or a proof that none exists for a stated class. \\
Multiplicative ergodic structure of scale cocycles & \cref{open:rg-multiplicative-ergodic} & A declared stationary family of blocking data over a measure-preserving base whose multiplier cocycle satisfies the hypotheses of a multiplicative ergodic theorem; or a proof that no such family exists for the exact interaction cocycle of the main manuscript. \\
Extensive relevance beyond the scalar Gaussian block & \cref{open:rg-hermite-extensive-relevance} & For a correlated, anisotropic, or gauge-mixed block, a declared transported family of laws and tangent norms together with a diagonalization of its score operator. \\
Nonlinear attraction at the Gaussian block law & \cref{open:rg-hermite-nonlinear-attraction} & A remainder estimate for the nonlinear block map in the score-tangent topology on a declared neighborhood of the Gaussian law. \\
Canonical continuous beta & \cref{open:rg-continuous-beta} & A canonical interpolation of a discrete blocking sequence with a proof of its uniqueness in a declared class; or a proof that the continuous beta is the same for every member of that class. \\
Two-index limits and universality & \cref{open:rg-two-index-limits} & A triangular family of finite agent networks at $c_*$ with projective consistency at fixed depth, uniform control of the coarse error, a diagonal scaling of depth with size, and a proof of convergence of a declared scheme-invariant observable; or an impossibility theorem for a stated class. \\
Primitive hierarchical identification & \cref{open:rg-scalarized-attraction} & For a stated family of graphs, a hierarchical identification under which sum-only aggregation yields one fixed primitive coupling matrix, with a Perron vector that factorizes for the bi-additive reading; or a proof that none exists. \\
\multicolumn{3}{@{}l}{\textbf{Part B-III: the lattice cluster}} \\
Typing of the lattice sites & \cref{open:rg-lattice-typing} & A declared reading proved: either the lattice law realized as the link marginal of a normalized joint with belief and model channels on a countable agent set at $c_*$ satisfying the countable-index projective and specification hypotheses, or a base discretization with a connection whose transports are the raw links and whose holonomies are the plaquette characters, with the path channel a scale map over a declared base map. \\
Limits beyond the admitted high-temperature construction & \cref{open:rg-spatial-generality} & New parameter domains with the corresponding tightness, interaction, and scaling estimates; positive Fisher floors uniform in depth; joint volume and depth estimates for the labeled experiment; and an existence and regularity theorem for a normalized law whose curvature action is computed from itself. \\
Quenched evidence density for general records & \cref{open:rg-pressure-quenched-record} & An every-exterior comparison and a pressure identity for site-dependent shifted fields of both signs in $(-1,1)$, together with an ergodic theorem for the record-dependent evidence, giving almost-sure record-independent limits of the evidence density, the posterior label rates, and the specific free energy. \\
\end{longtable}
}

\section{Numerical index}
\label{sec:rg-thermo-numerical-provenance}

\Cref{tab:base-numerical} indexes the computations recorded for this supplement. The first three rows
are the statements of this supplement whose status is NUMERICAL; their instances are deterministic,
and the implementation that computed them has been retired, so the values were not recomputed. The
remaining rows are finite exact comparisons for the lattice part, made in exact
rational arithmetic without random sampling. They corroborate established statements, which rest on
their proofs, and they supply no estimate as the lattice grows.

{\footnotesize
\renewcommand{\arraystretch}{1.12}
\begin{longtable}{@{}>{\raggedright\arraybackslash}p{0.16\textwidth}
>{\raggedright\arraybackslash}p{0.19\textwidth}
>{\raggedright\arraybackslash}p{0.25\textwidth}
>{\raggedright\arraybackslash}p{0.18\textwidth}
>{\raggedright\arraybackslash}p{0.11\textwidth}@{}}
\caption{Numerical index of Supplement B.}
\label{tab:base-numerical}\\
\toprule
\textbf{Claim} & \textbf{Check} & \textbf{What is tested} & \textbf{Control} & \textbf{Seed} \\
\midrule
\endfirsthead
\toprule
\textbf{Claim} & \textbf{Check} & \textbf{What is tested} & \textbf{Control} & \textbf{Seed} \\
\midrule
\endhead
\midrule
\multicolumn{5}{r@{}}{\emph{continued on the next page}} \\
\endfoot
\bottomrule
\endlastfoot
\cref{rem:rg-compatibility-computation} & \texttt{CHK-RG-\allowbreak RESCALING-\allowbreak COMPATIBILITY} & Blocking by the full ratio against blocking in declared stages, for one declared rescaling family on nested cycles of six categorical agents. & A flat-link control, the 24 orderings of the edges, and gauge covariance of the composite step. & no randomness \\
\cref{rem:rg-passive-blocking-computation} & \texttt{CHK-RG-\allowbreak PASSIVE-\allowbreak BLOCKING} & Iteration of the projected frozen-attention map at ratio two on a homogeneous six-cycle and ratio three on a homogeneous nine-cycle, and the pairwise block of its fixed structures. & The two ratios against each other, and the regeneration rule of the next row; no other control. & no randomness \\
\cref{rem:rg-regenerated-attention-computation} & \texttt{CHK-RG-\allowbreak REGENERATED-\allowbreak ATTENTION} & The declared row-regeneration rule on the same instances, and the pairwise block of its fixed structures. & The frozen-attention map from the same initial coupling, and a readout of equal cardinality that carries no gauge charge. & no randomness \\
\cref{thm:rg-thermo-native-cfour-pressure,prop:rg-thermo-native-ctwo-metric} & \texttt{CHK-RG-\allowbreak THERMO-\allowbreak NATIVE-\allowbreak JETS} & Exhaustive summation over the 256 sign configurations of the side-two square at $(b,t)=(1/4,1/20)$, with one internal plaquette: every raw-factor and physical-partition derivative through total degree three. & An independent normalized-score Gram comparison of the Fisher matrix at derivative degree zero. & no randomness \\
\cref{prop:rg-thermo-native-ctwo-metric} & \texttt{CHK-RG-\allowbreak THERMO-\allowbreak SINGULAR-\allowbreak FISHER} & The side-one square at $(b,t)=(1/3,1/20)$: two signs, native partition one, a null plaquette coordinate, and the field Fisher derivatives. & None separate. & no randomness \\
\cref{prop:rg-thermo-label-representatives} & \texttt{CHK-RG-\allowbreak THERMO-\allowbreak POSTERIOR-\allowbreak PUSH} & Two side-one branches with global factor normalization and a binary sensor of mean $k_\gamma b$, $k_\gamma\in\{1/2,1/3\}$, on one edge: posterior derivatives through total degree three against an independent two-coordinate Taylor-quotient calculation, and the induced pushed posterior and evidence derivatives. & The native partition is identically one, so a posterior with a nonconstant native partition divisor is not exercised. & no randomness \\
\end{longtable}
}
